\documentclass[../../../include/open-logic-section]{subfiles}

\begin{document}

\olfileid{sth}{choice}{vitali}
\olsection{ضميمه: د ويتالي پاراډوکس}

د باناخ--تارسکي د جوړښت د منلو يا نه منلو د سمې ارزونې
دپاره بايد د هغه \emph{ثبوت} وګورو۔ له بده مرغه، دا تر هغه
ډېر جبر غواړي چې دلته يې بيانولے شو۔ خو د لاندې پايلې
په رڼا کښې څو لنډې خبرې کولے شو چې ښايي د باناخ--تارسکي
ثبوت روښانه کړي۔\footnote{د ډېر بشپړ بحث دپاره
\cite{Weston2003} يا \cite{Wagon2016} وګورئ۔}

\begin{thm}[د ويتالي پاراډوکس (په $\ZFC$ کښې)]\ollabel{vitaliparadox}
هره دايره د شمېر وړ شمېر ټوټو ته وېشل کېدے شي؛ بيا دا
ټوټې (په ګرځولو او لېږدولو سره) داسې يوځاے کېدے شي
چې د هماغې دايرې دوه نقلونه ترې جوړ شي۔
\end{thm}

د ويتالي پاراډوکس ثبوت د باناخ--تارسکي د پاراډوکس تر
ثبوت ډېر اسانه دے۔ موږ دا د «ويتالي پاراډوکس» بولو، ځکه
چې د ويتالي د \citeyear{Vitali1905} د نه اندازېدونکي سټ
له جوړښت څخه راځي۔ خو د دواړو ثبوتونو سټي اړخونه ډېر
ورته دي۔ بنسټيز توپير يوازې دا دے چې باناخ--تارسکي
\emph{متناهي} وېش کاروي، او د ويتالي پاراډوکس
\emph{د شمېر وړ نامتناهي} وېش۔ د \citet{Weston2003}
په وينا، د ويتالي پاراډوکس د باناخ--تارسکي هومره حيرانوونکے
نه دے، خو ښيي چې هندسي پاراډوکسونه په «ساده» حالاتو
کښې هم پېښېدے شي۔

د ويتالي پاراډوکس له دوه-بُعدي شکل، يعنې دايرې، سره
اړيکه لري۔ نو په مستوي $\Real^2$ کښې کار کوو۔
$\rotationsgroup$ دې د مبدأ شاوخوا د ټکو د هغو
(د ساعت د ستنې په لوري) ګرځونو سټ وي چې د راډيان
زاويې يې په $[0,2\pi)$ کښې \emph{نسبي} عددونه دي۔
د $\rotationsgroup$ په اړه څو جبري حقيقتونه وړاندې کوو
(که د پايلې وينا اوس نه پوهېږئ، ثبوت به يې معنا روښانه کړي):

\begin{lem}\ollabel{rotationsgroupabelian}
$\rotationsgroup$ د تابعو د ترکيب له مخې يو تبديلېدونکے {ګروپ} جوړوي۔
\end{lem}

\begin{proof}
د $0_{\rotationsgroup}$ نښه د $0$ راډيانو د ګرځون دپاره
کاروو؛ دا د
$\rotationsgroup$ دپاره د هويت غړے دے، ځکه چې
$\comp{0_{\rotationsgroup}}{\rho} = \comp{\rho}{0_{\rotationsgroup}} =
\rho$ د هر $\rho \in \rotationsgroup$ دپاره۔

هر غړے معکوس لري۔ که $\rho \in \rotationsgroup$ په
$r$ راډيانو ګرځوي، نو $\rho^{-1} \in \rotationsgroup$
په $2\pi - r$ راډيانو ګرځوي، ځکه
$\rho \circ \rho^{-1} = 0_\rotationsgroup$۔

ترکيب تجميعي دے: $\comp{\rho}{(\comp{\sigma}{\tau})} =
\comp{(\comp{\rho}{\sigma})}{\tau}$ د هر $\rho, \sigma, \tau \in
\rotationsgroup$ دپاره۔

ترکيب تبديلېدونکے دے: $\comp{\rho}{\sigma} =
\comp{\sigma}{\rho}$ د هر $\rho, \sigma \in \rotationsgroup$ دپاره۔
\end{proof}
د چاپ شوي تعريف له مخې د نسبي راډياني زاويې دا سټ
د ترکيب لاندې تړلے نه دے: د څلورو راډيانو دوه ګرځونونه
د اته راډيانو ګرځون جوړوي، او په ورکړې وقفې کښې د هغه
همارزه زاويه غير نسبي ده۔ د پورته معکوس دعوه هم په همدې
تعريف کښې نه سمېږي؛ چاپي فورمولونه نه دي بدل شوي
(OLSTH-029)۔

په حقيقت کښې موږ خپل ګروپ $\rotationsgroup$ په دوو برخو
وېشلے شو، او بيا له هرې برخې ټول ګروپ بېرته ترلاسه کولے شو:

\begin{lem}\ollabel{disjointgroup}
$\rotationsgroup$ په دوو بېلو سټونو،
$\rotationsgroup_{1}$ او $\rotationsgroup_{2}$، وېشل کېدے شي
چې هر يو يې د $\rotationsgroup$ دپاره بنسټ دے۔
\end{lem}

\begin{proof}
په $\rotationsgroup_{1}$ کښې دې د $[0, \pi)$ د نسبي
راډياني زاويې ګرځونونه وي؛ او
$\rotationsgroup_{2} = \rotationsgroup
\setminus \rotationsgroup_1$ ونيسئ۔ د لومړني جبر له مخې،
$\Setabs{\comp{\rho}{\rho}}{\rho \in \rotationsgroup_1} =
\rotationsgroup$۔ د $\rotationsgroup_2$ دپاره هم ورته
پايله ترلاسه کېدے شي۔
\end{proof}

موږ به د ګروپونو په اړه همدا خبره د \olref{vitaliparadox}
د ثابتولو دپاره وکاروو۔ $\onesphere$ دې واحده دايره وي،
يعنې د مستوي له مبدأ څخه عيناً~$1$ واحد لرې د ټکو سټ،
يعنې $\Setabs{\tuple{r,s} \in \Real^2}{\sqrt{r^2+s^2}=1}$۔
د $\onesphere$ د وېشلو دپاره په $\onesphere$ دا اړيکه
په پام کښې نيسو:
\[
	r \sim s \emph{ هله او يوازې هله چې }(\exists \rho \in \rotationsgroup)\rho(r) = s.
\]
يعنې د $\onesphere$ دوه ټکي د دې اړيکې له مخې سره
تړلي دي که او يوازې که يو له بله د مبدأ شاوخوا په
نسبي ارزښت لرونکي ګرځون ترلاسه کېدے شي۔ نو تمه کېږي چې:

\begin{lem}
$\sim$ د همارزۍ اړيکه ده۔
\end{lem}

\begin{proof}
د \olref{rotationsgroupabelian} په کارولو ساده ده۔
\end{proof}

اوس د انتخاب اصل کاروو، څو داسې سټ $C$ واخلو چې د
$\onesphere$ د $\sim$ اړيکې د هرې همارزۍ طبقې څخه
عين يو غړے ولري۔ يعنې د همارزۍ د طبقو پر سټ د انتخاب
تابعه $f$ په پام کښې نيسو،\footnote{ځکه چې
$\rotationsgroup$ !!{enumerable} دے، د $E$ هر
!!{element} هم !!{enumerable} دے۔ او ځکه چې
$\onesphere$ !!{nonenumerable} دے، د
\olref{vitalicover} او
\olref[card-arithmetic][simp]{kappaunionkappasize}
له مخې $E$ هم !!{nonenumerable} دے۔ نو دا د
\emph{نه}شمېر وړ انتخاب کارونه ده۔}
\[
	E = \Setabs{\equivrep{r}{\sim}}{r \in \onesphere},
\]
او $C = \ran{f}$ ونيسئ۔ د هر ګرځون
$\rho \in \rotationsgroup$ دپاره، په $\funimage{\rho}{C}$
کښې هغه ټکي دي چې د $\rho$ ګرځون د $C$ پر هر ټکي
له کارولو ترلاسه کېږي۔ لاندې دوه پايلې ښيي چې دا سټونه
دايره بشپړه او بې له ګډو غړو پوښي:

\begin{lem}\ollabel{vitalicover}
$\onesphere = \bigcup_{\rho \in \rotationsgroup} \funimage{\rho}{C}$.
\end{lem}

\begin{proof}
$s \in \onesphere$ ثابت ونيسئ؛ داسې $r \in C$ شته چې
$r \in \equivrep{s}{\sim}$، يعنې $r \sim s$، يعنې
$\rho(r) = s$ د کوم $\rho \in \rotationsgroup$ دپاره۔
\end{proof}

\begin{lem}\ollabel{vitalinooverlap}
که $\rho_1 \neq \rho_2$، نو
$\funimage{\rho_1}{C} \cap \funimage{\rho_2}{C} = \emptyset$۔
\end{lem}

\begin{proof}
فرض کړئ $s \in \funimage{\rho_{1}}{C} \cap \funimage{\rho_{2}}{C}$۔
نو $s = \rho_{1}(r_{1}) = \rho_{2}(r_{2})$ د ځينو
$r_{1}, r_{2} \in C$ دپاره۔ له دې
$\rho^{-1}_2(\rho_1(r_1)) = r_2$، او
$\comp{\rho_1}{\rho^{-1}_2} \in \rotationsgroup$؛ نو
$r_{1} \sim r_{2}$۔ ځکه $r_1 = r_2$، چې $C$ د
$\sim$ د هرې همارزۍ طبقې څخه عين يو غړے ټاکي۔
نو $s = \rho_1(r_1) = \rho_2(r_1)$، او له دې
$\rho_1 = \rho_2$۔
\end{proof}

اوس خپل مخکيني جبري حقيقتونه په دايره کاروو:

\begin{lem}\ollabel{pseudobanachtarski}
$\onesphere$ په دوو بېلو سټونو، $D_{1}$ او $D_{2}$،
وېشل کېدے شي؛ $D_{1}$ بيا د شمېر وړ شمېر سټونو ته
وېشل کېدے شي چې په ګرځولو يې د $\onesphere$ يو نقل
جوړېږي (او د $D_{2}$ دپاره هم همداسې)۔
\end{lem}

\begin{proof}
د \olref{disjointgroup} څخه $\rotationsgroup_{1}$ او
$\rotationsgroup_{2}$ وکاروئ، او ونيسئ:
\begin{align*}
	D_{1} &= \bigcup_{\rho \in \rotationsgroup_1} \funimage{\rho}{C} & 
	D_{2} &= \bigcup_{\rho \in \rotationsgroup_2} \funimage{\rho}{C}
\end{align*}
دا د \olref{vitalicover} له مخې د $\onesphere$ وېش دے،
او $D_1$ او $D_2$ د \olref{vitalinooverlap} له مخې بېل دي۔
د جوړښت له مخې $D_1$ د شمېر وړ شمېر سټونو ته وېشل
کېدے شي: $\funimage{\rho}{C}$ د هر
$\rho \in R_1$ دپاره۔ دا سټونه په ګرځولو سره د
$\onesphere$ يو نقل جوړوي، ځکه چې د
\olref{disjointgroup} او \olref{vitalicover} له مخې
$\onesphere = \bigcup_{\rho \in
\rotationsgroup}\funimage{\rho}{C} = \bigcup_{\rho \in
\rotationsgroup_1}\funimage{(\comp{\rho}{\rho})}{C}$۔
د $D_2$ دپاره هم همدا استدلال چلېږي۔ \end{proof}\noindent
له دې سمدستي د ويتالي پاراډوکس راځي۔ ځکه د
$\onesphere$ څخه د $\onesphere$ \emph{دوه} نقلونه
جوړوو: په د شمېر وړ شمېر ټوټو (بېلابېلو
$\funimage{\rho}{C}$) يې وېشو، بيا يې په صُلب
حرکتونو خوځوو: د $D_1$ هره ټوټه يوازې
ګرځوو، او د $D_2$ هره ټوټه لومړے لېږدوو او بيا ګرځوو۔

د ثبوت تګلاره بيا وګورو۔ لومړے مو په مستوي کښې د
ګرځونو د ګروپ ځينې جبري حقيقتونه وکارول۔ په دې ګروپ
مو $\onesphere$ د همارزۍ په طبقو ووېشه۔ بيا مو د
انتخاب په کارولو له هرې طبقې يو غړے وټاکه او
«پاراډوکس» ته ورسېدو۔

د باناخ--تارسکي د ثبوت دپاره هم عين همدا تګلاره
کاروو۔ لوی توپير دا دے چې د باناخ--تارسکي د ثبوت
جبري حقيقتونه د ويتالي د پاراډوکس تر جبري حقيقتونو
ډېر پېچلي دي۔ خو دا جبري حقيقتونه له انتخاب سره
اړيکه نه لري۔ په لنډه توګه يې بيانوو۔

د باناخ--تارسکي د ثبوت دپاره د \olref{disjointgroup}
په څېر پايله لومړے ثابتوو: هر \emph{ازاد ګروپ} په
څلورو ټوټو وېشل کېدے شي، چې په تصوري ډول يې
«خوځولو» سره د ټول ګروپ دوه نقلونه ترلاسه کوو۔\footnote{د
\emph{څلورو} ټوټو بسيا کېدل د \cite{Robinson1947}
پايله ده۔ د نوي ثبوت دپاره \citet[Theorem 5.2]{Wagon2016}
وګورئ۔ د ګروپ د ټوټو د «خوځولو» په بيان کښې د
\citet[p.~3]{Weston2003} پيروي کوو۔} بيا ښيو چې
د $\Real^3$ د مبدأ شاوخوا دوه ځانګړي ګرځونونه د
ګرځونو يو ازاد ګروپ، $F$، جوړولے شي۔\footnote{وګورئ
\citet[Theorem 2.1]{Wagon2016}۔} (تر اوسه انتخاب نشته۔)
اوس د کرې د سطحې دوه ټکي «ورته» بولو که او يوازې
که يو يې په~$F$ کښې په کوم ګرځون له بله ترلاسه کېږي۔
بيا له «ورته» ټکو د همارزۍ د هرې طبقې څخه د عين يو
ټکي د ټاکلو دپاره \emph{انتخاب کاروو}۔ د
\olref{pseudobanachtarski} په څېر د $F$ وېش د کرې
پر سطحه کاروو او سطحه په څلورو ټوټو وېشو؛ دا ټوټې
«خوځوو» او د کرې د سطحې دوه نقلونه ترې جوړوو۔
له دې \citep{Hausdorff1914} ثابتېږي:

\begin{thm}[د هاوسدورف پاراډوکس (په $\ZFC$ کښې)]
د هرې کرې سطحه په متناهي شمېر ټوټو وېشل کېدے شي؛
دا ټوټې (په ګرځولو او لېږدولو سره) بيا داسې يوځاے
کېدے شي چې د هماغې کرې دوه بېل نقلونه ترې جوړ شي۔
\end{thm}

د باناخ--تارسکي بشپړې قضيې دپاره (چې يوازې د کرې
سطحه نه، بلکې دننه برخه يې هم رانغاړي) څو نورې جبري
لارې پکار دي۔ خو په رښتيا دا د جبري کار وروستۍ
پراختيا ده۔ ځکه وېسټن ليکي:
\begin{quote}
  [\ldots] د ازادو ګروپونو پايله د باناخ--تارسکي
  پاراډوکس د ثبوت \emph{بنسټيز ګام} دے۔ له دې نظره
  باناخ--تارسکي پاراډوکس تر دې چې د $\Real^3$ په اړه
  وينا وي، د هغو ګرځونو [او لېږدونو د ګروپ] د پېچلتيا
  په اړه وينا ده چې په $\Real^3$ کښې کار کوي۔
  \cite[p.~16]{Weston2003}
\end{quote}
يعنې متناهي وېش (لکه باناخ--تارسکي) او د شمېر وړ
نامتناهي وېش (لکه د ويتالي پاراډوکس) د دوه-بُعدي
او درې-بُعدي فضا د ګروپونو په ځينو حقيقتونو پورې تړلي دي۔

منو چې دا وروستۍ خبره د \olref[banach]{sec} د پای ټوکه
لږه خرابوي۔ ځکه چې ليکل په دوه بُعدونو کښې وي، نو که
«باناخ--تارسکي» په «باناخ--تارسکي باناخ--تارسکي»
بدلوو، بايد په \emph{د شمېر وړ نامتناهي} شمېر ټوټو
ووېشل شي۔ د ټوکې د سمولو دپاره بايد په درې بُعدونو
کښې وليکل شي۔ دا د لوستونکي د تمرين دپاره پرېږدو۔

يوه وروستۍ خبره۔ په \olref[banach]{sec} کښې مو وويل
چې د کرې ترلاسه شوې «ټوټې» ټولې \emph{د اندازې وړ}
نه شي کېدے؛ دا ناليدل کېدونکې «نامتناهي خورې ورې
ټولګې» دي۔ د \olref{pseudobanachtarski} په ترلاسه
کولو کښې د انتخاب کارونه هم همداسې ده۔ دا خبره تشريح
ارزي۔

بيا لږ مخينه بيانوو (خو دا \emph{يوازې} يوه لنډه نقشه
ده؛ د \emph{اندازې} په اړه د درسي کتاب برخه هم کتلے شئ)۔
د سټ $X$ دپاره د اندازې تعريف دا دے: ارزښت
$\mu(E) \in \Real$ د هر $E$ دپاره په کوم
«$\sigma$-الجبر» کښې، چې پر $X$ وي، ټاکل کېږي۔
دلته تفصيل اړين نه دے،
خو تابعه $\mu$ بايد د شمېر وړ جمعيت اصل ومني: د
بېلو سټونو د شمېر وړ اتحاد اندازه د هغو د بېلو اندازو
مجموعه ده، يعنې
$\mu(\bigcup_{n < \omega} X_n) = \sum_{n < \omega}\mu(X_n)$
کله چې $X_n$ سټونه دوه په دوه بېل وي۔ يو سټ «نه
اندازېدونکے» بلل دا دي چې ورته مناسبه اندازه نه شي
ټاکل کېدے۔ اوس خپل مخکينے $\rotationsgroup$ کاروو:

\begin{cor}[ويتالي]
$\mu$ دې داسې اندازه وي چې $\mu(\onesphere) = 1$،
او $\mu(X) = \mu(Y)$ وي که $X$ او $Y$ هم‌نهشت وي۔
بيا $\funimage{\rho}{C}$ د
هر $\rho \in \rotationsgroup$ دپاره نه اندازېدونکے دے۔
\end{cor}

\begin{proof}
د تناقض له لارې فرض کړئ چې داسې نه ده۔ نو
$\mu(\funimage{\sigma}{C}) = r$ د کوم
$\sigma \in \rotationsgroup$ او کوم $r \in \Real$
دپاره ونيسئ۔ د هر $\rho \in C$ دپاره
$\funimage{\rho}{C}$ او $\funimage{\sigma}{C}$
هم‌نهشت دي؛ نو $\mu(\funimage{\rho}{C}) = r$
د هر $\rho \in C$ دپاره۔ د \olref{vitalicover} او
\olref{vitalinooverlap} له مخې
$\onesphere = \bigcup_{\rho \in \rotationsgroup}\funimage{\rho}{C}$
د دوه په دوه بېلو سټونو د شمېر وړ اتحاد دے۔ نو د شمېر
وړ جمعيت له مخې $\mu(\onesphere) = 1$ د هر
$\funimage{\rho}{C}$ د اندازو مجموعه ده، يعنې
\[
	1 = \mu(\onesphere) = \sum_{\rho \in \rotationsgroup}\mu(\funimage{\rho}{C}) = \sum_{\rho \in \rotationsgroup}r
\]
خو که $r = 0$ وي، نو $\sum_{\rho \in \rotationsgroup}r = 0$؛
او که $r > 0$ وي، نو
$\sum_{\rho \in \rotationsgroup}r = \infty$۔
\end{proof}
په چاپي ثبوت کښې د ګرځون نښه دوه ځله د ټاکل شويو
ټکو په سټ کښې محدودېږي، حال دا چې د انځور تابع
ګرځون غواړي۔ چاپي کمّي کوونکے نه دے بدل شوے
(OLSTH-030)۔

\end{document}